\section{Introduction}
\label{sec:intro}

In modern deep learning, the ability to build generalist models that can solve multiple distinct tasks remains a primary objective.\footnote{Our code, converged task expert checkpoints, and evaluation scripts are publicly available at \href{https://github.com/anonymous-researcher/bc-router}{https://github.com/anonymous-researcher/bc-router}.} A highly successful paradigm is \emph{parameter-space model merging}, which fuses multiple specialized models (experts) fine-tuned from the same pre-trained base model into a single unified network without requiring additional training data or expensive computational resources \cite{task_arithmetic, ties_merging}. Standard approaches such as Task Arithmetic \cite{task_arithmetic} linearly interpolate the task-specific parameter differences (task vectors). More advanced static merging protocols, such as TIES-Merging \cite{ties_merging} and DARE \cite{dare}, attempt to resolve spatial and sign interference by filtering low-magnitude weights and resolving parameter conflicts coordinate-wise.

While static merging provides a computationally cheap, training-free model fusion, it imposes a rigid parameter state that must process all input samples identically. Consequently, static merging suffers from severe multi-task representation conflict; optimizing the global static scaling coefficient to favor one task inevitably degrades performance on another, a phenomenon known as representation collapse. To overcome this, the field has recently shifted toward \emph{dynamic model merging}, where the parameter-space merging coefficients are dynamically routed on a per-input or per-batch basis.

A prominent and highly visible representative of this dynamic merging trend is \emph{Quantum Wavefunction Superposition Merging} (QWS-Merge) \cite{qws_merge}. QWS-Merge challenges the static parameter assumption by modeling task-specific experts as quantum eigenstates in a parameter Hilbert space. By projecting the input representation onto a unit sphere and formulating routing coefficients using cosine phase-overlap equations (resembling wave-like phase interference), QWS-Merge claims to achieve state-of-the-art (SOTA) multi-task routing. Crucially, the authors of QWS-Merge report that a standard classical Linear Router baseline catastrophically collapses on high-conflict tasks such as SVHN.

In this work, we adopt a critical, methodological perspective and apply Occam's razor to these claims. As methodologists, we argue that bad experimental design, under-tuned baselines, and a lack of fully converged expert models lead to false progress in the deep learning community. In particular, we expose that when task-specific experts are trained to true convergence, a close mathematical and structural audit of QWS-Merge deconstructs the necessity of its quantum wavefunction superposition metaphor. Specifically, we identify **four major methodological insights**:
\begin{enumerate}
    \item \textbf{Paradigm Distinction vs.\ Performance Similarity:} Under a fair and converged evaluation pipeline, QWS-Merge achieves a joint homogeneous accuracy of $83.97 \pm 0.53$\%, while the unsupervised test-time adaptation (TTA) baseline AdaMerging achieves $89.30 \pm 0.00$\%. However, we highlight that comparing them purely on accuracy conflates two fundamentally different operational paradigms: AdaMerging dynamically optimizes merging coefficients on-the-fly during inference via active gradient descent to minimize prediction entropy, which introduces significant latency, backward-pass overhead, and susceptibility to temporal stream noise. In contrast, offline-calibrated dynamic routing (like QWS-Merge or our proposed routers) optimizes its routing head once on a tiny calibration set prior to deployment, running as a pure, lightweight forward pass with zero test-time latency or optimization overhead.
    \item \textbf{Deconstructing Classical Failures via Proper Regularization:} The unregularized classical Linear Router baseline collapses on SVHN ($74.00 \pm 16.14$\%) and exhibits extreme instability due to low-data overfitting on the tiny calibration set. When we apply standard L2 regularization during calibration, our simple, classical \textbf{Linear Router (Reg)} completely resolves this collapse, boosting SVHN accuracy to \textbf{91.73} $\pm$ \textbf{3.71}\% (outperforming QWS-Merge by \textbf{+12.00\%}) and increasing the joint homogeneous accuracy to $82.80 \pm 4.85$\%. This proves that the reported failure of classical linear routers is a pure baseline tuning artifact.
    \item \textbf{Resolving the Softmax Zero-Sum competitive Bottleneck:} Standard classical routers employ Softmax routing, which forces tasks to compete in a zero-sum game for a shared "routing budget" during mixed-batch calibration. This causes BL-Router's SVHN collapse ($31.73 \pm 16.03$\%). We introduce the Softmax-free \textbf{BSigmoid-Router}, which replaces Softmax with independent Sigmoids, completely resolving this competitive bottleneck to achieve a stable joint homogeneous accuracy of $83.73 \pm 1.93$\% and heterogeneous stream accuracy of \textbf{83.96} $\pm$ \textbf{2.27}\% ($B=1$, outperforming QWS-Merge's $83.29 \pm 0.36$\%).
    \item \textbf{The Layer-wise Specialization Confounder \& Wave Regularization:} QWS-Merge has a massive layer-wise capacity advantage over global baselines. When we control for this confounder using our regularized \textbf{GLS-Router}, we demonstrate that simple, classical layer-specific routing can achieve comparable multi-task behavior, but exhibits extreme sensitivity to the calibration set across seeds (producing standard deviations as high as \textbf{24.30\%} on SVHN). This reveals that QWS-Merge's true value lies in its wave equations acting as a stable, structural regularizer.
\end{enumerate}

To evaluate these methods fairly and establish high-fidelity results, we fine-tune specialized experts on a Vision Transformer (\texttt{vit\_tiny\_patch16\_224}) backbone to true convergence across four distinct, high-conflict vision tasks (MNIST, FashionMNIST, CIFAR-10, SVHN). Our empirical findings yield a stark and compelling deconstruction of the quantum model merging narrative, demonstrating that simple classical routers, when properly regularized or formulated without the Softmax zero-sum bottleneck, match or outperform complex wave superposition metaphors.

By deconstructing these mechanisms, we apply Occam's razor to demystify dynamic model merging. Our work serves as a critical call to action for the machine learning community: we must prioritize rigorous, fair evaluation protocols, control for key structural confounders, and recognize that exposing flaws and identifying the true drivers of performance—such as proper regularization—is just as valuable as proposing complex new architectures.